\documentclass[11pt]{article}
\usepackage{fontspec,amsmath,amssymb,amsthm,hyperref}
\usepackage[margin=1in]{geometry}
\setlength{\emergencystretch}{3em}
\title{Reference drift moments uniform over smoothing scales}
\author{LangevinLean mathematical exposition}
\date{}
\begin{document}\maketitle
\noindent Source module: \texttt{ProximalLangevin/SplittingMoments.lean}. Compile this proof document with XeLaTeX or LuaLaTeX. The accompanying HTML places the complete Lean declarations beside every proof.


For \(F=f+e\), assume \(F\) is \(m\)-strongly convex, \(f,e\) differentiable, \(\nabla f\) is \(L\)-Lipschitz, and \(\|\nabla e\|\le K\). The goal is a Gibbs drift bound using \(L,m,K\), without an inverse smoothing scale.


\section{The envelope preserves a global value Lipschitz constant}


Let \(p=P_\lambda(y)\). As a competitor at \(x\), use \(p+x-y\), so its quadratic penalty equals that of \(p\) at \(y\). The \(K\)-Lipschitz inequality for \(f\) changes the function value by at most \(K\|x-y\|\). Thus \(f^\lambda(x)\le f^\lambda(y)+K\|x-y\|\), and reversing \(x,y\) proves the same global Lipschitz constant.



\noindent Lean declarations: \texttt{moreau\_lipschitz\_of\_lipschitz}.

\section{Find a true minimum and bound the smooth force there}


Strong convexity gives quadratic coercivity. Minimize \(F\) on a sufficiently large closed ball, with radius \(2\|\nabla F(0)\|/m+1\); the quadratic lower bound proves the resulting point minimizes globally. At a global minimum \(z\), the local-extremum derivative vanishes. Hence \(\nabla f(z)=-\nabla e(z)\) and \(\|\nabla f(z)\|\le K\).



\noindent Lean declarations: \texttt{gradient\_eq\_zero\_at\_minimum}, \texttt{exists\_strongConvex\_minimum}, \texttt{gradient\_smoothPart\_at\_minimum}.

\section{A bounded gradient remains bounded in L\ensuremath{{}^{2}}}


The norm of the derivative of a \(K\)-Lipschitz scalar function is at most \(K\); the gradient is its Riesz representative and has the same norm. For a pointwise bounded drift \(\|b(x)\|\le K\), almost everywhere representative equality gives \(\|b(Y)\|_{L^2}\le K\) on a probability space.



\noindent Lean declarations: \texttt{gradient\_norm\_le\_of\_lipschitz}, \texttt{lpDrift\_norm\_le\_of\_bounded}.

\section{Integrate a sharp quadratic bound about the mode}


Lipschitz continuity and the mode bound give \(\|\nabla f(x)\|\le L\|x-z\|+K\), so \(\|\nabla f(x)\|^2\le2L^2\|x-z\|^2+2K^2\). Sharp Gibbs concentration supplies \(\mathbb E_\pi\|x-z\|^2\le d/m\), yielding

\[\mathbb E_\pi\|\nabla f\|^2\le2L^2d/m+2K^2.\]

For any \(Y\) with actual Gibbs law \(\pi\), the pushforward integral formula transfers this estimate and taking square roots gives \(\|\nabla f(Y)\|_2\le\sqrt{2L^2d/m+2K^2}\).



\noindent Lean declarations: \texttt{integral\_sq\_smooth\_gradient\_le}, \texttt{lpDrift\_smooth\_gradient\_norm\_le}.

\end{document}
